\chapter{总结与展望}
\label{chap:conclusion}

\section{全文工作总结}
\label{sec:conclusion-summary}

本文研究医学影像任务变化与训练信息受限条件下的持续学习。全文以 MedCL 平台设计和持续分割基准组织评价问题，分别研究弱监督持续分割、类别扩展下的联邦持续分类，以及跨器官和跨模态持续配准。

第三章给出 MedCL 平台设计，拟接收本地训练完成的阶段模型或预测，在固定测试集上评分，生成阶段--任务结果及分客户端分析。平台部分形成需求、架构、接口和测试方案，尚未完成端到端实现与系统验证。

在平台设计之后，第三章通过持续分割基准研究域变化、类别扩展和跨器官任务变化，在统一病例划分、网络、任务顺序与训练配置下比较代表性方法。六中心序列中，ER 的最终平均 Dice 为 $0.750\pm0.021$，高于 EWC 的 $0.672\pm0.021$，而二者 E-FWT 接近。这说明最终性能、历史保持、当前学习和前向泛化具有不同含义，需要结合参数增长与历史数据访问条件共同评价。既有基准为平台接口和结果展示提供实证参照，其实验独立于平台完成。

第四章直接承接六中心分割的数据与测试设置，将训练监督改为固定模拟涂鸦。以静态弱监督方法 ZScribbleSeg 为基础，ZSDERpp 将全局一致性、空间先验、有限图像--涂鸦回放、参考特征约束与回放弱监督优化结合。在六中心域增量序列中，其最终平均 Dice 为 $0.701$，BWTR 为 $-0.152$，均高于增强涂鸦顺序训练，但 E-FWT 未达到最高值。结果支持当前配置下的最终性能与历史保持改善，也表明稀疏监督覆盖、当前域适应和未来域泛化需要分别分析。

第五章研究类别持续扩展、原始数据分布于多个客户端且历史原始样本不回放的联邦分类。FedSubMerge 递归构建并融合客户端主梯度子空间，通过正交补投影减少新更新对跨客户端历史方向的一阶干扰；FedSubMerge-AD 按层选择相关客户端，以协调历史覆盖和当前可塑性。PathMNIST 的 $\alpha=0.3$ 设置下，FedSubMerge-AD 的 ACC 为 $93.42\%$，较 FOT 提高 $7.22$ 个百分点；统一版本在七个设置中均取得高于自适应版本的 BWT。实验采用同步类别任务流、已知测试任务身份及任务特定分类头。模型和子空间仍保存历史摘要，子空间交换本身不提供形式化隐私保证。

第六章提出 SAMCL，结合有限历史图像对回放、一阶元持续学习与锐度感知优化。脑部任务的最终 Dice 从直接顺序训练的 $0.348$ 提高到 $0.697$，后向变化从 $-0.395$ 提高到 $-0.037$。加入锐度感知优化后，所考察的域内和跨任务泛化结果有所改善。SAMCL 与 MER 在不同任务上各有优势，配准精度与历史保持需分别分析。

三项方法保存的历史信息不同：ZSDERpp 使用图像--涂鸦对和参考特征，FedSubMerge 使用梯度子空间摘要，SAMCL 使用固定--移动图像对。它们在不同训练条件下处理新任务学习与旧任务保持的矛盾。共同评价关注阶段性能变化，各任务保留自己的输出定义和指标。

\section{研究局限性}
\label{sec:conclusion-limitations}

MedCL 尚未完成端到端实现，其运行效率、并发处理、安全隔离和使用效果仍待验证。持续分割基准采用预先划分的任务、固定骨干和公开数据，未覆盖连续漂移、中心重访和未预先标明任务切换时点的长期数据流。

弱监督分割使用由密集掩膜生成的模拟涂鸦，未测量真实标注时间、标注者差异和漏标错标影响；模型选择仍依赖密集验证标签。定量实验限于六中心域增量，类别扩展和跨器官适配尚未验证。ZSDERpp 的回放容量为 64 幅，第三章默认容量为 32 幅，跨章比较同时涉及监督、回放规模和方法组成的变化。

联邦分类采用同步任务流、预先给定的任务切换时点和已知测试任务身份，服务器接收并融合客户端梯度子空间。梯度投影只作用于选定共享参数，未投影参数、模型漂移和截断误差仍可能影响历史性能。客户端由公开数据构造，通信核算未覆盖真实网络时延、故障和多服务器调度；低秩梯度摘要也不提供形式化隐私保证。

持续配准依赖有限历史原始图像对。四项任务同时涉及器官、模态、配准关系和损失配置变化，尚未分离这些因素的影响。实验只采用固定顺序和 LK-UNet 骨干，尚未检验不同任务顺序、其他形变模型及长期数据流中的任务切换识别。

三项方法分别研究分割、分类和配准，并未由同一个模型联合学习三类任务。现有实验主要采用回顾性公开数据，真实多中心纵向验证和临床工作流程中的使用效果仍待研究。

\section{未来研究展望}
\label{sec:conclusion-future}

\subsection{医学多模态模型的持续更新}

未来可将本文的持续学习研究扩展到医学多模态基础模型，考察影像和报告持续到达时，模型如何更新表示、适应不同中心与下游任务，并修订经审核的医学知识。图\ref{fig:future-medical-mllm}展示了这一设想，研究重点是保持既有的跨模态关联和历史任务能力。

\begin{figure}[htbp]
\centering
\resizebox{\textwidth}{!}{% Report and scan illustrations feed a shared model, followed by a task timeline.
\begingroup
\linespread{1}\selectfont
\definecolor{modelteal}{HTML}{24777B}
\definecolor{modelink}{HTML}{37474F}
\begin{tikzpicture}[
  font=\thesisfigurefont,
  box/.style={draw=modelink!70,rounded corners=2pt,line width=.65pt,align=center,inner sep=0pt},
  flow/.style={-{Latex[length=2mm]},draw=modelink,line width=.75pt,shorten <=1pt,shorten >=1pt},
  labeltext/.style={align=center,inner sep=0pt},
  task/.style={box,minimum width=3.75cm,minimum height=2.65cm,fill=black!1}
]
% Editable report sheets: visual placeholders only, without patient information.
\begin{scope}[shift={(1.05,-.05)}]
  \draw[draw=modelink!35,fill=black!3,rounded corners=1pt] (.14,.08) rectangle (1.69,1.63);
  \draw[draw=modelink!45,fill=white,rounded corners=1pt] (.07,.04) rectangle (1.62,1.59);
  \path[draw=modelink,line width=.7pt,fill=white] (0,0) -- (1.55,0) -- (1.55,1.24) -- (1.25,1.55) -- (0,1.55) -- cycle;
  \draw[draw=modelink,line width=.6pt] (1.25,1.55) -- (1.25,1.24) -- (1.55,1.24);
  \draw[draw=modelteal,line width=1.4pt] (.21,1.23) -- (.51,1.23) (.36,1.08) -- (.36,1.38);
  \draw[draw=modelink!70,line width=.7pt] (.7,1.28) -- (1.06,1.28) (.7,1.12) -- (1.12,1.12);
  \draw[draw=modelink!45,line width=.55pt] (.22,.87) -- (1.29,.87) (.22,.65) -- (1.29,.65) (.22,.43) -- (1.03,.43);
  \draw[draw=modelteal!60,line width=.55pt] (.22,.2) -- (.76,.2);
\end{scope}
\node[labeltext,anchor=west,font=\thesisfigurefont\bfseries] at (2.95,.98) {医学文本};
\node[labeltext,anchor=west,text=black!65] at (2.95,.45) {报告\enspace /\enspace 病史};
\node[inner sep=0pt] at (8.65,.7) {\includegraphics[height=1.5cm]{ch07/assets/example-000.png}};
\node[labeltext,anchor=west,font=\thesisfigurefont\bfseries] at (10.0,.98) {医学影像};
\node[labeltext,anchor=west,text=black!65] at (10.0,.45) {持续到达的数据};
\node[box,minimum width=2.8cm,minimum height=.6cm,fill=white] (textenc) at (3.025,-.95) {文本编码器};
\node[box,minimum width=2.8cm,minimum height=.6cm,fill=white] (imageenc) at (10.075,-.95) {图像编码器};
\draw[flow] (3.025,-.12) -- (textenc.north);
\draw[flow] (10.075,-.12) -- (imageenc.north);
\node[box,draw=modelteal,line width=1pt,fill=modelteal!5,minimum width=13.1cm,minimum height=1.55cm] (model) at (6.55,-2.48) {};
\node[labeltext,text=modelteal,font=\thesisfigurefont\bfseries] at (6.55,-2.02) {医学多模态基础模型};
\node[labeltext] at (6.55,-2.5) {影像\quad 文本\quad 结构\quad 知识};
\node[labeltext,text=black!65] at (6.55,-2.96) {跨模态融合\enspace /\enspace 共享表征\enspace /\enspace 版本记录};
\draw[flow] (textenc.south) -- (model.north -| textenc.south);
\draw[flow] (imageenc.south) -- (model.north -| imageenc.south);
\node[task] (task1) at (1.875,-5.48) {};
\node[task] (task2) at (6.55,-5.48) {};
\node[task] (task3) at (11.225,-5.48) {};
\node[labeltext,font=\thesisfigurefont\bfseries] at (1.875,-4.65) {任务一\\胸部 X 线诊断};
\node[labeltext,font=\thesisfigurefont\bfseries] at (6.55,-4.65) {任务二\\消化内镜识别};
\node[labeltext,font=\thesisfigurefont\bfseries] at (11.225,-4.65) {任务三\\皮肤镜分类};
\node[inner sep=0pt] at (1.875,-5.95) {\includegraphics[height=1.35cm]{ch07/assets/example-001.png}};
\node[inner sep=0pt] at (6.55,-5.95) {\includegraphics[height=1.35cm]{ch07/assets/example-002.png}};
\node[inner sep=0pt] at (11.225,-5.95) {\includegraphics[height=1.35cm]{ch07/assets/example-003.png}};
\draw[flow] (model.south -| task1.north) -- (task1.north);
\draw[flow] (task1.east) -- (task2.west);
\draw[flow] (task2.east) -- (task3.west);
\node[labeltext,text=black!65] at (8.45,-3.72) {新模态、新任务、新中心与新知识};
\draw[flow,dashed,draw=modelteal!70] (3.75,-4.02) -- (13.1,-4.02);
\node[labeltext,text=black!65] at (6.55,-7.18) {顺序适配\quad $\longrightarrow$\quad 阶段评价};
\draw[draw=black!20,line width=.55pt] (0,-7.59) -- (13.1,-7.59);
\node[labeltext] at (6.55,-8.04)
  {\textbf{稳定性}\enspace 既有能力保持\qquad $\longleftrightarrow$\qquad \textbf{可塑性}\enspace 新增信息吸收};
\end{tikzpicture}
\endgroup
}
\caption[医学多模态基础模型持续学习的未来研究框架]{医学多模态基础模型持续学习的未来研究框架。持续到达的影像与文本经多模态基础模型形成共享表示，并依次适配不同影像模态、中心和下游任务。未来研究从数据与表征、任务与场景、知识与行为三个层次展开，分别对应持续预训练、参数高效持续适配和持续知识编辑。}
\label{fig:future-medical-mllm}
\end{figure}

更新对象可分别位于影像与文本表示、任务适配模块，以及模型使用的知识或规则。相应地，未来研究分为持续预训练、参数高效持续适配和持续知识编辑三个方向，并分别确定可访问历史信息、测试任务与资源需求。

\begin{figure}[htbp]
\centering
\resizebox{\textwidth}{!}{% Three research directions surround a shared objective, not a serial algorithm.
\begingroup
\linespread{1}\selectfont
\definecolor{routegreen}{HTML}{3D7166}
\definecolor{routegold}{HTML}{AB7E32}
\definecolor{routerust}{HTML}{A65D48}
\begin{tikzpicture}[
  font=\thesisfigurefont,
  direction/.style={circle,minimum size=3.75cm,draw,line width=.85pt,inner sep=0pt},
  labeltext/.style={align=center,inner sep=0pt},
  spoke/.style={-{Latex[length=2mm]},draw=black!45,line width=.75pt,shorten <=2pt,shorten >=2pt}
]
\node[direction,draw=routegreen,fill=routegreen!5] (data) at (1.95,0) {};
\node[direction,draw=routegold,fill=routegold!6] (task) at (11.15,0) {};
\node[direction,draw=routerust,fill=routerust!5] (knowledge) at (6.55,-4.2) {};
\node[labeltext,text=routegreen] at (1.95,.7) {数据与表征};
\node[labeltext,font=\thesisfigurefont\bfseries] at (1.95,.08) {持续预训练};
\node[labeltext,text=black!65] at (1.95,-.6) {新数据\quad 表征保持};
\node[labeltext,text=routegold] at (11.15,.90) {任务与场景};
\node[labeltext,font=\thesisfigurefont\bfseries] at (11.15,.02) {参数高效\\持续适配};
\node[labeltext,text=black!65] at (11.15,-.85) {新场景\quad 参数高效};
\node[labeltext,text=routerust] at (6.55,-3.5) {知识与行为};
\node[labeltext,font=\thesisfigurefont\bfseries] at (6.55,-4.12) {持续知识编辑};
\node[labeltext,text=black!65] at (6.55,-4.8) {新知识\quad 能力保持};
\node[circle,minimum size=3.05cm,draw=black!65,fill=black!3,line width=.9pt,inner sep=0pt] (goal) at (6.55,-.6) {};
\node[labeltext,font=\thesisfigurefont\bfseries] at (6.55,-.37) {医学多模态\\持续学习};
\node[labeltext,text=black!65,font=\sffamily\heiti\fontsize{10bp}{12bp}\selectfont] at (6.55,-1.12) {稳定性\,$\leftrightarrow$\,可塑性};
\draw[spoke] (goal) -- (data);
\draw[spoke] (goal) -- (task);
\draw[spoke] (goal) -- (knowledge);
\end{tikzpicture}
\endgroup
}
\caption[医学多模态基础模型持续学习的三层研究路线]{医学多模态基础模型持续学习的三层研究路线。持续预训练更新影像与文本表示，参数高效持续适配处理中心和任务变化，持续知识编辑修订经审核的知识与输出规则。MedCL 的后续实现拟记录模型版本，并比较各版本的历史能力与新增能力。}
\label{fig:future-three-level-roadmap}
\end{figure}

\subsection{数据与表征层：持续预训练与表征保持}

持续预训练需要在吸收新影像和新报告时，保留已有器官结构、疾病相关视觉特征和影像--文本对应关系。可在历史访问受限条件下，比较少量历史样本、旧模型输出和压缩表征对这些能力的保持效果。

这一层次需要建立带有时间与来源记录的数据流，区分采集时间、报告版本、患者划分和来源中心。对于影像与报告不完全配对或部分模态缺失的情况，应分别评价单模态能力和跨模态匹配，避免以样本数量增长替代数据质量分析。少量代表性历史样本、旧模型输出或压缩表征可作为知识保持的候选载体，但其存储成本、信息保真和共享限制需要在相同条件下比较。

评价包括固定历史任务、新任务迁移和跨模态匹配。分割、分类和配准用于检查结构与语义表示，影像检索和图文对应用于检查跨模态关联。各更新版本与未更新模型及相同新增数据下的直接训练模型比较，并报告训练步数、显存与历史信息存储量。

\subsection{任务与场景层：参数高效持续适配}

新中心的设备、患者构成、标注规范和下游用途可能与既有数据不同。后续可研究低秩适配、提示或小型任务模块，在限制可训练参数数量的条件下持续适应新场景，并结合历史特征保持、子空间保护或有限回放。

当数据未标明任务切换时点，且旧中心可能再次出现时，需要判断何时更新、何时复用已有适配模块，以及何时增加新模块。评价同时记录任务或模块选择错误及模块数量增长。联邦环境还需区分本地适应与跨客户端共享，研究哪些更新适合在客户端之间交换。

在固定基础模型、数据划分和训练预算的条件下，比较完整微调、独立适配与持续适配，并分别记录可训练参数数量、计算量和存储开销。除新中心性能外，评价旧中心保持、未知中心泛化和模块选择错误；联邦设置还报告通信量与客户端参与情况。

\subsection{知识与行为层：持续知识编辑与能力保持}

经审核的医学知识、术语或输出规则发生变化时，可研究针对性的连续编辑，更新指定内容并保持无关知识及已有任务能力。重点比较连续多次编辑的累积影响，以及影像概念与文字描述之间的对应关系是否发生非预期变化。

对于医学多模态模型，编辑目标既可能涉及文本知识，也可能涉及影像概念与文字描述之间的对应关系。需要区分更新外部检索资料、修改提示或规则以及直接修改模型参数三种方式，分别说明修改作用于哪一层、何时生效以及能否撤回。新增知识条目应具有可核验来源和版本；涉及医学判断的更新还需要相应专业审核，不能仅以模型生成的新回答作为修改正确性的依据。

评价可围绕编辑目标的正确性、同义提问和相关影像上的泛化、无关知识与旧任务保持，以及连续编辑后的累积干扰展开。对于证据不足的输入，还应检查模型是否保留适当的不确定性表达和人工复核入口。本文尚未开展知识编辑实验，后续需先建立明确的编辑集合、未编辑对照和历史能力测试，再判断局部更新是否带来可接受的能力变化。

\subsection{MedCL 的实现与验证}

MedCL 后续可接收外部更新后的模型或预测，在固定测试集上比较不同版本。持续预训练、持续适配和知识编辑分别报告表示保持、新旧场景表现，以及指定修改的正确性和影响范围。

平台实现优先完成基准选择、阶段配置、模型或预测提交、独立评分和报告生成。验证内容包括样本匹配、任务身份与分类头选择、配准坐标约定、缺失阶段处理，以及模型推理和上传文件的安全检查。联邦评价按固定客户端测试分组评分。各项后续算法研究独立训练，其模型或预测按规定格式提交评测。
